\documentclass[11pt,a4paper]{article}
% Common preamble for the hint sheets (\input from each file)
\usepackage[utf8]{inputenc}
\usepackage[T1]{fontenc}
\usepackage{lmodern}
\usepackage{amsmath,amssymb,amsthm,mathtools}
\usepackage[a4paper,margin=2.3cm]{geometry}
\usepackage{enumitem}
\usepackage{xcolor}
\usepackage[hidelinks]{hyperref}
\renewcommand{\leq}{\leqslant}
\renewcommand{\le}{\leqslant}
\renewcommand{\geq}{\geqslant}
\renewcommand{\ge}{\geqslant}
\setlength{\parindent}{0pt}
\setlength{\parskip}{0.35em}
% \hintstitle{Sheet title}{Class work / Homework}
\newcommand{\hintstitle}[2]{\begin{center}{\LARGE\bfseries Hints: #1\par}\vspace{0.5em}{\large #2\par}\end{center}
\noindent\textit{The hints to each problem are cumulative. Read one, try again, and open the next only if you are still stuck.}\par\medskip\hrule\medskip}
\newcommand{\problem}[1]{\par\bigskip{\large\bfseries Problem #1}\par\nopagebreak\smallskip}
\newcommand{\subpart}[1]{\par\smallskip\textbf{(#1)}\par\nopagebreak}
\newcommand{\hint}[1]{\par\smallskip\noindent\textcolor{black!70}{\textbf{Hint #1.}}\enspace}
\newcommand{\note}[1]{\par\smallskip\noindent\textit{Note.} #1\par}
\providecommand{\lt}{<}
\providecommand{\gt}{>}
\providecommand{\tightlist}{\setlength{\itemsep}{0pt}\setlength{\parskip}{0pt}}
\emergencystretch=3em

\begin{document}
\hintstitle{Integrals -- 3. Differentiation under the integral sign and Frullani integrals}{Homework}

\problem{1}
\hint{1}
Let \(F(\alpha)=\int_0^\infty e^{-\alpha x}\frac{\sin\beta x}xdx\). Then \(F'(\alpha)=-\int_0^\infty e^{-\alpha x}\sin\beta x\,dx=-\frac{\beta}{\alpha^2+\beta^2}\) (uniform convergence for \(\alpha\ge\alpha_0\gt0\)), and \(F(\alpha)\to0\) as \(\alpha\to\infty\).
\hint{2}
Hence \(F(\alpha)=\arctan\frac\beta\alpha\). Letting \(\alpha\to0^+\) (justified as in the class sheet, Problem 4) gives \(\int_0^\infty\frac{\sin\beta x}xdx=\frac\pi2\operatorname{sgn}\beta\).

\problem{2}
\hint{1}
\(F'(\alpha)=\int_0^1\frac{dx}{(1+\alpha^2x^2)\sqrt{1-x^2}}\), and \(F(0)=0\). Substitute \(x=\sin\theta\).
\hint{2}
\(\int_0^{\pi/2}\frac{d\theta}{1+\alpha^2\sin^2\theta}=\frac{\pi}{2\sqrt{1+\alpha^2}}\) (use \(t=\tan\theta\)). Integrating, \(F(\alpha)=\frac\pi2\ln\left(\alpha+\sqrt{1+\alpha^2}\right)=\frac\pi2\operatorname{arsinh}\alpha\).

\problem{3}
\hint{1}
For \(\alpha\gt0\): \(F'(\alpha)=\int_0^\infty\frac{2\alpha x^2}{(1+\alpha^2x^2)(1+x^2)}dx\). Decompose into partial fractions in \(x^2\): \(\frac{x^2}{(1+\alpha^2x^2)(1+x^2)}=\frac1{\alpha^2-1}\left(\frac1{1+x^2}-\frac1{1+\alpha^2x^2}\right)\).
\hint{2}
This gives \(F'(\alpha)=\frac{\pi}{1+\alpha}\), and \(F(0)=0\), so \(F(\alpha)=\pi\ln(1+\alpha)\) for \(\alpha\ge0\) (and \(\pi\ln(1+\lvert\alpha\rvert)\) in general). Check at \(\alpha=1\): \(x=\tan\theta\) gives \(-2\int_0^{\pi/2}\ln\cos\theta\,d\theta=\pi\ln2\).

\problem{4}
\hint{1}
Look for \(f(x,t)\) on \([0,1]\) whose \(\partial_tf\) blows up near the corner \((x,t)=(0,0)\). Try \(f(x,t)=\frac{xt^3}{(x^2+t^2)^2}\), \(f(x,0)=0\).
\hint{2}
Directly, \(F(t)=\int_0^1f\,dx=\frac{t}{2(1+t^2)}\), so \(F'(0)=\frac12\). But for each fixed \(x\), \(\partial_tf(x,0)=\lim_{t\to0}\frac{xt^2}{(x^2+t^2)^2}=0\), so the formal answer is \(0\). The hypothesis that fails: \(\partial_tf\) is not continuous on \([0,1]\times[-\delta,\delta]\) (at \(x=t\) it equals \(\frac1{4t}\), unbounded). Another standard example: \(\int_0^\infty\frac{\sin\alpha x}xdx\), whose formal derivative \(\int_0^\infty\cos\alpha x\,dx\) diverges.

\problem{5}
\hint{1}
This is Frullani with \(f(u)=e^{-u}\), \(a=1\), \(b=\cosh t\): the answer is \(\ln\cosh t\).
\hint{2}
The two pieces \(\int_0^\infty\frac{e^{-x}}x\) and \(\int_0^\infty\frac{e^{-x\cosh t}}x\) diverge separately, so no splitting is legitimate on \((0,\infty)\). On \([\varepsilon,\infty)\) both converge and their difference is \(\int_\varepsilon^{\varepsilon\cosh t}\frac{e^{-u}}udu\), which tends to \(\ln\cosh t\); this limit is just the improper integral of the whole integrand, so it does not depend on how you split.

\problem{6}
\hint{1}
Write the integrand as an integral: \(\frac{x-1}{\ln x}=\int_0^1x^y\,dy\) and \(\frac{e^{-ax}-e^{-bx}}x=\int_a^be^{-yx}\,dy\). Swap the order (justify by Fubini/Tonelli: the integrands are nonnegative, or absolutely integrable).
\hint{2}
First: \(\int_0^1\frac{dy}{y+1}=\ln2\). Second: \(\int_0^\infty e^{-yx}\cos cx\,dx=\frac y{y^2+c^2}\), so the answer is \(\int_a^b\frac{y\,dy}{y^2+c^2}=\frac12\ln\frac{b^2+c^2}{a^2+c^2}\).

\problem{7}
\hint{1}
\(\partial_a I=\int_0^{\pi/2}\frac{2a\cos^2x}{a^2\cos^2x+b^2\sin^2x}dx\). Divide by \(\cos^2x\) and set \(t=\tan x\): \(\int_0^\infty\frac{2a\,dt}{(a^2+b^2t^2)(1+t^2)}\).
\hint{2}
Partial fractions give \(\partial_aI=\frac{\pi}{a+b}\), so \(I=\pi\ln(a+b)+C(b)\). At \(a=b\) the integral is \(\pi\ln b\), hence \(C=-\pi\ln2\) and \(I=\pi\ln\frac{a+b}2\).

\problem{8}
\hint{1}
Rewrite each power with base \(e\): \(2^{\log x}=e^{\log2\cdot\log x}=x^{\log2}\); \(x^{1/\log x}=e^{\log x/\log x}=e\); \(2015^x=e^{x\log2015}\); and \(e^{e^{2016x}+6048x}=e^{e^{2016x}}\left(e^{2016x}\right)^3\).
\hint{2}
Answers: \(\frac{2^{1+\log2}}{1+\log2}\); \(e\cdot e=e^2\) (the integrand is the constant \(e\)); \(\frac{2015^x}{\log2015}+C\); with \(u=e^{2016x}\), \(\frac1{2016}\int u^2e^u\,du=\frac{e^{u}(u^2-2u+2)}{2016}+C\).

\problem{9}
\hint{1}
Frullani needs only that \(f(0)\) and \(f(\infty)\) exist, not that \(f(\infty)=0\). Here \(f=\arctan\), \(f(0)=0\), \(f(\infty)=\frac\pi2\) (for \(a,b\gt0\)).
\hint{2}
The answer is \(\left(0-\frac\pi2\right)\ln\frac ba=\frac\pi2\ln\frac ab\). Check: differentiating in \(a\) gives \(\int_0^\infty\frac{dx}{1+a^2x^2}=\frac{\pi}{2a}\).

\problem{10}
\hint{1}
By parts with \(u=\sin^2x\), \(dv=\frac{dx}{x^2}\): the boundary terms vanish and you get \(\int_0^\infty\frac{\sin2x}xdx=\frac\pi2\).
\hint{2}
With \(F(\alpha)=\int_0^\infty\frac{\sin^2\alpha x}{x^2}dx\): \(F'(\alpha)=\int_0^\infty\frac{\sin2\alpha x}xdx=\frac\pi2\) for \(\alpha\gt0\) and \(F(0)=0\), so \(F(\alpha)=\frac{\pi\alpha}2\) (uniform convergence of the derivative integral holds on \([\alpha_0,\infty)\); use continuity at \(0\)).
\hint{3}
Parseval: for \(f=\mathbf 1_{[-1,1]}\), \(\hat f(\omega)=\int f(x)e^{-i\omega x}dx=\frac{2\sin\omega}\omega\), and \(\int\lvert\hat f\rvert^2d\omega=2\pi\int\lvert f\rvert^2dx=4\pi\). So \(\int_{-\infty}^\infty\frac{\sin^2\omega}{\omega^2}d\omega=\pi\) and the requested integral is \(\frac\pi2\).
\end{document}
